\documentclass[a4paper]{article}
\usepackage{MMVII}

\newcommand{\TheCmd}{OriPoseSelecAllPAir}

\begin{document}

\CmdTitle
%\maketitle

\tableofcontents


 %=================================================================================

\SecDescr

This command is used in global relative orientation, it computes a set of pairs  and triplets
\footnote{probably, later, the triplet computation will be separated} for which relative orientation
will be computed.
Generally it will not be called \emph{individually} but rather \emph{implicitelly} inside \FileSty{OriPoseEstimRelGlob}.
However  there is two possible motivation for \emph{explicit} call to this command :

\begin{itemize}
   \item {\bf didactic}, to understand step by step the global orientation process
   \item {\bf fine tuning } when the default parametrization used by \FileSty{OriPoseEstimRelGlob}
         is not adapted.
\end{itemize}



 %=================================================================================

\section{Mandatory Parameters and "optional mandatory" parameter}

 % ---------------------------------------------- 
\subsection{Mandatory Parameters}

The two mandatory parameters are the set of images and the destination 
folder that will contain all the result for relative orientation
(a subfolder of \FileSty{MMVII-PhgrProj/OriRel/}).


 % ---------------------------------------------- 
\SubSecTiePRelOr{The selection of pair is based on tie point.}

 %=================================================================================

\section{Methodology}

\subsection{Introduction}

The approach is quite basic for now and was developped to have a fast
prototype of global orientation, it will most certainly evolve.

It's difficult to formalize the problem rigourously as the selection
of pair/triplet  optimal according to criterion.  By the way, practically 
the requirement is to have sufficient set of pair/triplet, to allow
a good reconstruction, and not too many to limit computation time.

So the approach used is to have several heuristics corresponding to different criteria
and to accumulate their results.  These heuristics require some quality evaluation of the 
pairs, for now it's limited to the  number of tie-point (one aspect that will need to be improved).


We create a symetric graph $G_{tp}$ such that :

\begin{itemize}
   \item there is a vertice for eachf image;
   \item there is an edge between all image where the number $N$ of tie 
         points is over the parameter \FileSty{NbMinTieP};
    \item the weight of each edge is $\frac{1}{1+N}$ (any postive decreasing function will do).
\end{itemize}


 % ---------------------------------------------- 
\subsection{Pair selection and control parameter}

The pair selection is made of $3$ heuristiks :

\begin{itemize}
   \item  to assure connectivity a minimal spaning $T_0$ tree is added, 
          to have better connectivity the optimal the optimal spaning $T_1$ tree in $G_0=G_{tp}-T_0$ is also added,
          then $T_2$ in  $G_1=G_0-T1$ \dots the number of iteration is controlled by \FileSty{NbMSTree};

   \item   for each image  we select its $K$  \emph{best} neighboors, where $K$ is controled
           by \FileSty{NbKBest} ;

   \item   using the alphabetic order, let $I_1,..,I_N$  be the images, we add all the pair 
           $I_j,I_k$ such that $(j-k)\%N|<\delta$  where $\delta$ is controled by \FileSty{NbSuccLines};
           this heuristic is usefull when the acquisition has some time coherence, which is generally true.

\end{itemize}

Finnaly for reasonably small dataset, we can enforce to select all the pair with parameter \FileSty{AllPair}.
Let note $E_s$ the edges selected.


 % ---------------------------------------------- 
\subsection{Triplet selection }

Considerring all the $3$ length loops of graph $G_{tp}$,  we add a triplet
corresponding to this loop is :


\begin{itemize}
   \item the three edges of this loop belong tp $E_s$;
   \item two  edges of this loop belong to $E_s$, and the missing edge
         is not the worst of loop.
\end{itemize}


Next step, is to creat a new  command for triplet selection, that will be called after pair orientation.

 %=================================================================================

\section{Result}


If \FileSty{REL} is the second mandatory parameter,
a subfolder of \FileSty{MMVII-PhgrProj/OriRel/}REL/) is created.
It will  containt file for storing pair and triplet selected.


\end{document}






